\documentclass[10pt]{article}
\usepackage[T1]{fontenc}
\usepackage[utf8]{inputenc}
\usepackage[french]{babel}
\usepackage{amsmath,amssymb,mathrsfs,mathpazo}
\usepackage{geometry}
\geometry{a4paper,margin=2cm}
\usepackage{graphicx,array,tabularx,tabularray,colortbl,multirow,multicol,booktabs,enumitem,pifont,ulem,fancybox,diagbox,cancel}
\usepackage{pstricks,pstricks-add,pst-plot,pst-tree,pst-node,pst-text,pst-func,pst-math,pst-bspline,pst-3dplot,pst-eucl,pst-solides3d}
\let\pscalculate\undefined
\usepackage{pst-bezier}
\usepackage{tikz,pgfplots,tkz-tab}
\usetikzlibrary{arrows,arrows.meta,patterns,calc,positioning,trees,decorations.markings,shapes.geometric,shapes.misc,shapes.arrows,intersections,decorations.pathmorphing,decorations.pathreplacing,backgrounds,shadows,babel}
\pgfplotsset{compat=1.18}
\usepgfplotslibrary{fillbetween}
\usepackage[np]{numprint}
\usepackage{listings,makecell,pas-tableur,twoopt}
\let\barre\undefined
\usepackage[pstricks]{tablvar}
\definecolor{darkgreen}{RGB}{0,100,0}
\providecommand{\red}{\color{red}}
\providecommand{\blue}{\color{blue}}
\providecommand{\green}{\color{green}}
\providecommand{\black}{\color{black}}
\providecommand{\white}{\color{white}}
\providecommand{\gray}{\color{gray}}
\providecommand{\cyan}{\color{cyan}}
\providecommand{\magenta}{\color{magenta}}
\providecommand{\yellow}{\color{yellow}}
\providecommand{\violet}{\color{violet}}
\providecommand{\orange}{\color{orange}}
\providecommand{\euro}{\texteuro}
\providecommand{\EUR}{\texteuro}
\providecommand{\ieme}{\textsuperscript{e}}
\providecommand{\degres}{\ensuremath{^\circ}}
\providecommand{\degree}{\ensuremath{^\circ}}
\providecommand{\vv}[1]{\overrightarrow{#1}}
\providecommand{\Coord}[1]{#1}
\providecommand{\decosix}{}
\providecommand{\starredbullet}{\ensuremath{\bullet}}
\pagestyle{empty}
\setlength{\parindent}{0pt}
\usepackage{adjustbox}\DeclareUnicodeCharacter{2236}{:}\DeclareUnicodeCharacter{2190}{\ensuremath{\leftarrow}}
\providecommand{\R}{\mathbb{R}}
\providecommand{\N}{\mathbb{N}}
\providecommand{\Z}{\mathbb{Z}}
\providecommand{\Q}{\mathbb{Q}}
\providecommand{\D}{\mathbb{D}}
\providecommand{\C}{\mathbb{C}}
\providecommand{\e}{{\mathrm{e}}}
\providecommand{\vect}[1]{\overrightarrow{#1}}
\providecommand{\vectt}[1]{\overrightarrow{\mathrm{#1}}}
\providecommand{\vvt}[1]{\overrightarrow{\mathrm{#1}}}
\providecommand{\np}[2][]{{#2~#1}}
\providecommand{\vv}[1]{{\overrightarrow{#1}}}
\renewcommand{\d}{{\mathrm{d}}}
\providecommand{\degres}{{^\circ}}
\providecommand{\degre}{{^\circ}}
\providecommand{\blue}{\color{blue}}
\providecommand{\red}{\color{red}}
\providecommand{\enskip}{\quad}
\providecommand{\llbracket}{\lbrack\!\lbrack}
\providecommand{\rrbracket}{\rbrack\!\rbrack}
\providecommand{\euro}{{\text{€}}}
\providecommand{\no}{{\text{n°}}}
\providecommand{\interval}[2]{{[#1;#2]}}
\providecommand{\intervFO}[2]{{[#1;#2[}}
\providecommand{\dotfill}{\ldots}

\providecommand{\ds}{\displaystyle}
\providecommand{\textding}[1]{#1}
\providecommand{\Oij}{\(\mathrm{O};\vec{\imath},\vec{\jmath}\)}
\providecommand{\Oijk}{\(\mathrm{O};\vec{\imath},\vec{\jmath},\vec{k}\)}
\providecommand{\Ouv}{\(\mathrm{O};\vec{u},\vec{v}\)}
\providecommand{\intervFF}[2]{\ensuremath{[#1;#2]}}
\renewcommand{\euro}{\ifmmode\text{\texteuro}\else\texteuro\fi}
\DeclareUnicodeCharacter{00B0}{\ifmmode^{\circ}\else\textdegree\fi}
\DeclareUnicodeCharacter{20AC}{\ifmmode\text{\texteuro}\else\texteuro\fi}
% Énoncés : Éducation nationale. Transcription et corrigés : APMEP, Denis Vergès et contributeurs. Découpage et classement : CoopMaths, Sébastien Lozano, Éric Elter et contributeurs. Adaptation de la présentation : Maths Explorer.
% CC BY-SA : https://coopmaths.fr/www/about/
% https://coopmaths.fr/alea/static/eam/2025/tex/eam_technologique_2025_06_sujet0v1_1_cor.tex
\begin{document}
\noindent\textit{APMEP / CoopMaths. CC BY-SA. Présentation adaptée par Maths Explorer.}\par\medskip
\smallskip

\textbf{Question 1}

Le dimanche le temps passé à faire les devoirs est $\dfrac{25}{100} = \dfrac{25 \times 1}{25 \times 4} = \dfrac 14$ ; sur ce temps le temps consacré a l'exposé est de 80\,\% soir $\dfrac{80}{100} \times \dfrac 14 = \dfrac 14 \times \dfrac{80}{100}$ : réponse B

\textbf{Question 2}

Diminuer un prix de 50\,\% c'est le multiplier par $1 - \dfrac{50}{100} = 1 - 0,5 = 0,5$.

Pour retrouver le prix initial il faut doubler ce pris soit l'augmenter de $1 + 1 = 1 + \dfrac{100}{100}$, soit l'augmenter de 100\,\%.

\textbf{Question 3 }

Pour obtenir 200 à partir de 250, il suffit de multiplier par 200 et de diviser par 250, soit de le multiplier par $200 \times \dfrac{1}{250} = \dfrac{200}{250} = \dfrac{200 \times 4}{250 \times 4} = \dfrac{800}{\np{1000}} = 0,8$

\textbf{Question 4}

On a bien $\dfrac{10^{-5}}{10^{8}} = 10^{-5 - 8} = 10^{- 13}$.

\textbf{Question 5}

L'épaisseur d'une pile de 2 000 feuilles est égale à $\np{2000} \times 70 \times 10^{-3} = 2 \times 7 \times 10^{3 + 1 -3} = 14 \times 10^{1} = 140$~(mm) ou 14~(cm).

\textbf{Question 6}

Terre : $5,973 \times 10^{24}$~(kg) ; Mercure : $3,302 \times 10^{23}$~(kg) ;

Vénus : $\np{4,8685} \times 10^{24}$~(kg) ; Mars : $\np{6,4185} \times 10^{23}$~(kg).

La masse la plus grande est celle de la Terre.

\textbf{Question 7}

On a la somme : $x + 3x + x^2 = 4x + x^2$.

\textbf{Question 8}

La courbe $\mathcal{C}'$ est au dessous de la courbe $\mathcal{C}$ lorsque $x$ appartient à l'intervalle $[-2~;~-1]$ ou lorsque $x$ appartient à l'intervalle $[1~;~2]$.

\textbf{Question 9}

La courbe coupe l'axe des abscisses en deux points d'abscisses négatives.

\textbf{Question 10}

La fonction s'annule en $x = 2$ ce qui élimine B et D.

On a $f(0) > 0$, ce qui élimine C : reste A.

\textbf{Question 11}

On a $C = (1 + t)^2$, donc $C > 0$ et $\sqrt{C}$ existe : on a donc :

$1 + t = \sqrt{C}$ ou $1 + t = - \sqrt{C}$, d'où :

$t = \sqrt{C} - 1$ ou $t = - \sqrt{C} - 1$

\textbf{Question 12}

C'est l'année 2016.

\bigskip

\end{document}
